% PAGE 1 -- Executive summary for busy management. Three headings ONLY.
% Prose is yours to write. The comments under each heading are the points the
% analysis supports (see notes/tencent_issue.md for the numbers and sources).
\section*{Executive summary}

\subsection*{What is the issue?}
% - Prosus owns 22.66% of Tencent and equity-accounts it under IAS 28 (an associate,
%   not a subsidiary and not a financial asset at fair value).
% - Carrying value US$42.2bn vs quoted value US$127.6bn at 31 March 2026: an
%   US$85.5bn gap that is in no statement, equal to 1.6x group equity.
% - FY2026 profit US$11.65bn = US$7.26bn share of Tencent's profit (non-cash)
%   + US$4.7bn gain on selling 0.9% of Tencent (the gap crystallising) - US$0.3bn
%   from everything else. Cash dividends from associates: US$1.24bn.
% - The treatment rests on "significant influence", which Prosus does not list
%   as a critical judgement, with a stake that falls ~0.8 pp a year towards 20%.
% Citations: FY2026 AR note 9 pp.155-156; income statement p.140; note 3 p.145.

\subsection*{Why is it relevant?}
% - Any multiple on reported earnings or book equity is wrong: earnings are
%   Tencent's, not cash, and partly a realisation of old gains; equity is
%   understated by ~2.6x. Value on look-through NAV instead.
% - After a deal the acquirer books Tencent at fair value (IFRS 3): the disposal
%   gains disappear, impairment headroom goes from US$85bn to nil, earnings get
%   lower and more volatile. The acquirer's accounts will not look like Prosus's.
% - Tax: Prosus books no deferred tax on the gap (Dutch participation exemption,
%   Tencent not taxed at source). Whether that survives the deal structure is a
%   US$10-20bn question (20% of the gap = US$17bn, illustrative).
% - Income the buyer cannot direct: Tencent sets dividends (~1% yield) and
%   buy-backs; US$17bn of bonds sit against the quoted value, not the book value.

\subsection*{What should we do?}
% - Proceed on a NAV basis: Tencent at market, other listed stakes at market,
%   unlisted at fair value, operated businesses on their own cash flows. Drop P/E.
% - Tax-structuring opinion before price: keep the shares inside MIH Internet
%   Holdings B.V.; test CFC/minimum-tax rules at home; map the Naspers chain.
% - Model the acquirer's post-deal accounts: Tencent at ~US$128bn, impairment
%   sensitivity to a 20-30% Tencent drawdown, covenants and rating in LTV terms.
% - Decide now what happens at 20%: keep selling (cliff in ~3 years, one-off
%   remeasurement, then FVOCI with ~US$1bn of dividend income as "profit"), or stop.
% - Diligence: Tencent's own December accounts and the Q1 lag adjustments; the
%   board seat and any shareholder arrangements; post-year-end sales (US$675m).
